\documentclass[11pt]{article}
\usepackage[a4paper,margin=25mm]{geometry}
\usepackage{amsmath,amssymb,amsthm,booktabs}
\usepackage[british]{babel}
\usepackage{url,xcolor,graphicx}
\usepackage[round,authoryear]{natbib}
\usepackage[hidelinks]{hyperref}
\hypersetup{pdftitle={Constrained static hedging of a capped GBM claim},
pdfauthor={Amanjeet Singh},pdfsubject={Physical loss and funding constraints},
pdfkeywords={static hedging, mean-square loss, convex optimisation, KKT}}
\usepackage{pgfplots}
\usepgfplotslibrary{groupplots}
\pgfplotsset{compat=1.18}
\setlength{\bibhang}{1.5em}
\setlength{\bibsep}{2pt}
\setlength{\emergencystretch}{1.5em}
\renewcommand{\bibfont}{\small}
\renewcommand{\topfraction}{0.85}
\renewcommand{\textfraction}{0.1}
\renewcommand{\floatpagefraction}{0.8}
\newtheorem{theorem}{Theorem}[section]
\title{Constrained static hedging of a capped GBM claim}
\author{Amanjeet Singh}
\date{11 October 2026}
\begin{document}
\maketitle
\begin{abstract}
I study static hedging of a capped European claim under geometric Brownian
motion with position and funding limits. Risk-neutral prices determine the
initial premium and instrument costs, while the physical measure determines
mean-square discounted terminal error. Dividends are reinvested explicitly.
Closed truncated-lognormal moments yield a strictly convex quadratic loss
and a unique constrained hedge. Complete optimality proofs include the
zero-budget boundary, a finite active-set method, the distinction between
variance and mean-square objectives, and budget shadow prices. I compare
126 constrained hedges with unrestricted and cash-only baselines, a
same-constraint variance optimiser, a projection heuristic and an independent
general-purpose solver. Monte Carlo and separate residual-square integration
check the risk calculations. Adverse heuristic rankings, small numerical
violations and unresolved cancellation cases remain visible. Funding limits
have an exact risk penalty; additional budget has a locally interpretable
marginal value. The study provides a reproducible financial optimisation
example with rigorous proofs, using synthetic inputs and established hedging
foundations.
\end{abstract}
\noindent\begin{minipage}{\linewidth}
\small\emergencystretch=2em
\textit{Keywords:} static hedging; mean-square loss; physical measure;
convex optimisation; funding constraints; KKT conditions; shadow prices.\par
\smallskip\textit{JEL classification:} C61, C63, G13.\par
\end{minipage}
\section{Financial model and funded loss}
\label{sec:h-model}

I study how position and funding constraints alter static hedging of a
capped claim under a classical geometric Brownian motion (GBM) model.
Static quadratic hedging with pricing-measure cost constraints is
established \citep{leunglorig2016}; I specialise it to a two-instrument
basis and a fixed no-arbitrage premium. Fixed-capital mean-square loss
also appears in constrained dynamic hedging \citep{czichowskyschweizer2012},
but the existence proof here is finite-dimensional.
\footnote{The closest static antecedent was read in its identified author
preprint; its published successor metadata was verified. The source ledger
records that distinction and the limited sections actually consulted.}
The underlying GBM market remains dynamically complete. Residual risk
arises from the restricted instrument set and static positions. Physical
parameters are stipulated for synthetic comparisons.

Let $S,K,H,T,\sigma>0$, let $r,q,\mu$ be finite constants, and set
$J=K+H$, $D=e^{-rT}$ and $a=e^{(q-r)T}$. Under the physical measure,
\begin{equation}
 X=S\exp\{(\mu-q-\sigma^2/2)T+\sigma\sqrt T Z\},\qquad
 Z\sim N(0,1),\qquad Y=D\min\{(X-K)^+,H\}.
 \label{eq:h-physical-law}
\end{equation}
Here $\mu$ is the total-return drift. I price the contract under the
stipulated equivalent measure $Q$, whose ex-dividend drift is $r-q$.
For strike $A$, let
\begin{align}
 C_A&=S e^{-qT}\Phi(d_{1,A})-A D\Phi(d_{2,A}),\notag\\
 d_{1,A}&=\frac{\log(S/A)+(r-q+\sigma^2/2)T}{\sigma\sqrt T},
 \quad d_{2,A}=d_{1,A}-\sigma\sqrt T.
 \label{eq:h-fair-premiums}
\end{align}
These classical prices follow \citet{blackscholes1973} and
\citet{merton1973}. The received premium is $V_0=C_K-C_J>0$. Under $Q$, terminal stock has a strictly positive
lognormal density on $(0,\infty)$. The non-negative call and capped
payoffs are positive on sets of positive probability, so their finite
discounted expectations give $C_K>0$ and $V_0>0$. I allow only the
total-return stock, with dividends reinvested, and the vanilla call
at $K$. Their initial costs and discounted gains are
\begin{equation}
 p=(S,C_K)^\top,\qquad
 g=(aX-S,\ D(X-K)^+-C_K)^\top.
 \label{eq:h-instruments}
\end{equation}
For fixed holdings $\theta=(\theta_s,\theta_c)^\top$, the initial bank
balance is $V_0-p^\top\theta$. The self-financing portfolio's discounted
terminal wealth is $V_0+\theta^\top g$, so the discounted liability
residual and loss are
\begin{equation}
 R(\theta)=Y-V_0-\theta^\top g,\qquad
 L(\theta)=\mathbb E_P[R(\theta)^2]
 =c-2h^\top\theta+\theta^\top M\theta,
 \label{eq:h-risk}
\end{equation}
where $M=\mathbb E_P[gg^\top]$, $h=\mathbb E_P[g(Y-V_0)]$ and
$c=\mathbb E_P[(Y-V_0)^2]$. This is uncentred mean-square financial
loss, including squared bias, rather than residual variance alone.
It penalises discounted terminal surplus and shortfall symmetrically;
the pricing premium is fixed, not chosen to minimise this loss.
I impose the declared box and instrument-spend budget
\begin{equation}
 {\cal F}_B=\{\theta:0\leq\theta_s,\theta_c\leq1,\
 p^\top\theta\leq B\},\qquad B=\kappa V_0,\quad \kappa\geq0.
 \label{eq:h-feasible-set}
\end{equation}
For $\kappa\leq1$ the bank balance is non-negative; for $\kappa>1$
borrowing may occur. These position bounds are modelling choices,
not universal limits on optimal hedges.

\section{Convexity and constrained optimality}

The relevant convex optimality principles are established
\citep{boydvandenberghe2004}. The following specialisation includes the
zero-budget boundary, where strict feasibility fails.

\begin{theorem}[Existence, uniqueness and optimality]
\label{thm:h-constrained-optimum}
The matrix $M$ is positive definite and $\mathcal F_B$ is nonempty
and compact. The constrained hedge is unique. Its necessary and
sufficient conditions are feasibility, non-negative multipliers
$\ell,u\in\mathbb R^2$ and $\lambda\in\mathbb R$, and
\begin{align}
 2(M\theta-h)-\ell+u+\lambda p&=0,\label{eq:h-kkt-stationarity}\\
 \ell_i\theta_i=0,\quad u_i(\theta_i-1)&=0,\quad
 \lambda(p^\top\theta-B)=0.\label{eq:h-kkt-complementarity}
\end{align}
These conditions also hold at $B=0$, although strict feasibility
fails there. The optimum lies among at most sixteen candidates:
the unconstrained stationary point, stationary points on five
supporting lines and their ten pairwise intersections, retaining
only well-defined feasible candidates.
\end{theorem}

Appendix~\ref{app:h-proofs} gives full proofs and closed Gaussian
moment expressions for every objective coefficient. It also expresses
risk as an integral of three squared affine residuals. No finite hedge
in this instrument span replicates the capped claim exactly; apparent
zero risk from subtracting nearly equal floating-point moments is
therefore a numerical warning. A long call at $K$ and short call at
$J$ would replicate the claim, but that second call is unavailable here.

\section{Loss choice and the value of funding}
\subsection{Variance minimisation and mean-square hedging}

The financial objective also explains why variance-reduction controls
cannot simply be reused as optimal hedges. Put $m=\mathbb E_Pg$,
$e=\mathbb E_P(Y-V_0)$, $G=\operatorname{Cov}_P(g)$ and
$k=\operatorname{Cov}_P(g,Y-V_0)$. With
$\theta_v=G^{-1}k$, $\zeta=m^\top G^{-1}m$ and
$\eta=e-m^\top\theta_v$, the unrestricted mean-square relaxation is
\begin{equation}
 \theta_M=M^{-1}h=\theta_v+\frac{G^{-1}m\,\eta}{1+\zeta},
 \qquad L(\theta_v)-L(\theta_M)=\frac{\eta^2\zeta}{1+\zeta}.
 \label{eq:h-centred-risk-gap}
\end{equation}
Under $P=Q$, $m=e=0$ and the objectives coincide; different measures
permit, but do not force, a strict difference. This unrestricted
$L^2$ benchmark can contain short positions and is not automatically
admissible under the bounded-below wealth condition. Constraints add the
exact penalty
$L(\theta_B)-L(\theta_M)=(\theta_B-\theta_M)^\top
M(\theta_B-\theta_M)$.

\subsection{Budget shadow prices}

The optimal risk $v(B)=\min_{\mathcal F_B}L$ is non-increasing
and convex. Every budget multiplier yields the global supporting bound
\begin{equation}
 v(B')\geq v(B)-\lambda_B(B'-B),\qquad B'\geq0.
 \label{eq:h-budget-shadow}
\end{equation}
On a regular active-set interval, $v'(B)=-\lambda_B$ and $v$ is
quadratic or constant. At zero budget,
$v'(0+)=-2\max\{0,h_s/S,h_c/C_K\}$: the initial marginal
benefit of available budget is the largest positive loss-reduction/cost
ratio; unhelpful spending is not required. The budget here permits spending; it does not require the entire
amount to be invested.
\section{Registered comparisons and numerical evidence}
\label{sec:h-results}

I retain the seven designed contracts in Table~\ref{tab:h-contracts}.
The protocol was saved before outcomes, with physical total-return drifts
$\mu\in\{0.03,0.06,0.12\}$ and budget multiples
$\kappa\in\{0,0.25,0.5,1,2,4\}$. This gives 126 constrained
mean-square comparisons. Rates are annualised; maturities are in years.
The parameters are synthetic, with no licensed or empirical observations.
The experiment assumes a known physical law; it does not fit $\mu$.

\begin{table}[tb]\centering\small
\begin{tabular}{lrrrrrrr}\toprule
Contract & $S$ & $K$ & $H$ & $T$ & $r$ & $q$ & $\sigma$\\\midrule
Lower kink & 100 & 100 & 20 & 1 & 0.03 & 0.01 & 0.2 \\
Upper kink & 120 & 100 & 20 & 1 & 0.03 & 0.01 & 0.2 \\
Inside spread & 110 & 100 & 20 & 1 & 0.03 & 0.01 & 0.2 \\
Deep cap & 160 & 100 & 20 & 1 & 0.03 & 0.01 & 0.2 \\
Short dated & 100 & 100 & 20 & 0.05 & 0.03 & 0.01 & 0.06 \\
Negative rate/dividend & 105 & 100 & 20 & 1 & -0.015 & 0.045 & 0.25 \\
High volatility & 95 & 100 & 30 & 2 & 0.03 & 0.01 & 0.65 \\
\bottomrule\end{tabular}
\caption{Designed one-asset inputs. Instrument costs and the received premium use $Q$; terminal loss uses the specified physical drift. No input is a market estimate.}
\label{tab:h-contracts}\end{table}

The analytic unconstrained solution is the strongest two-instrument
relaxation, and cash-only financing is the no-hedge baseline.
I compare the constrained solution with a separately solved
centred-variance objective using the same constraints, a Euclidean
projection heuristic, and SciPy's general-purpose SLSQP optimiser.
The projection heuristic minimises distance from unrestricted
variance-optimal holdings; it does not minimise financial loss.
The exact call-spread reference uses an additional call and a short
position, so it identifies the importance of the instrument restriction
rather than superiority within a shared feasible set.

Closed truncated moments determine objective coefficients.
Independent adaptive normal integration, split at both payoff kinks,
checks 168 mean/product coefficients. Risk is reported primarily by
integrating the non-negative squared residual, with the expanded
moment quadratic retained alongside it. I do not replace negative
computed risks with zero or repair eigenvalues. The numerical mean-square
matrix condition number reaches $42,337$, despite exact positive
definiteness. Thus numerical representation remains a separate issue
from the convexity proof.

\begin{table}[tb]\centering\small
\begin{tabular}{lrrrr}\toprule
Contract & No hedge & Constrained & Unrestricted & Reduction (\%)\\\midrule
Lower kink & 68.4032 & 15.7205 & 13.5707 & 77.02 \\
Upper kink & 58.6954 & 24.9738 & 20.2983 & 57.45 \\
Inside spread & 71.9973 & 22.9439 & 19.0941 & 68.13 \\
Deep cap & 6.06651 & 5.23772 & 4.218 & 13.66 \\
Short dated & 0.770884 & Unresolved & Unresolved & n/a \\
Negative rate/dividend & 85.0105 & 24.4634 & 20.1739 & 71.22 \\
High volatility & 151.748 & 77.0938 & 57.1103 & 49.20 \\
\bottomrule\end{tabular}
\caption{Population-risk diagnostics at $\mu=0.06$ and $B=V_0$. Losses are squared discounted monetary units. The percentage compares the constrained mean-square hedge with cash-only financing. Unrestricted holdings relax the box and budget. The short-dated hedge risks lie below a heuristic arithmetic floor and support neither a resolved risk ratio nor a zero-risk assertion.}
\label{tab:h-risk-comparison}\end{table}

No registered optimiser failure or quadrature warning occurred.
Maximum scaled primal violation is $1.88\times10^{-14}$; maximum
negative multiplier magnitude is $2.84\times10^{-14}$.
The largest stationarity and complementary-slackness residuals are
$4.55\times10^{-13}$ and $2.76\times10^{-15}$ respectively.
These are floating-point diagnostics, not exact feasibility or
admissibility certificates. SLSQP's residual-square risks differ from
the active-set values by between $-2.27\times10^{-13}$ and
$1.59\times10^{-11}$, with maximum position difference
$2.34\times10^{-7}$. The small negative comparison remains visible.
Equation~\eqref{eq:h-kkt-risk-bound} states the exact prerequisites
needed to translate residuals into a risk bound.

Among the 117 constrained rows above the heuristic floor, reductions
against the no-hedge baseline range from zero to $81.5868$ per cent.
The largest same-constraint variance-objective penalty is $0.191175$
squared monetary units, or $0.5982$ per cent of the mean-square
optimum, for the negative-rate/dividend contract at $\mu=0.12$,
$\kappa=0.5$. The distinction is mathematically substantive but
quantitatively modest in this registered family. Euclidean projection
has a largest penalty of $612.244$, illustrating why geometric
closeness to an unconstrained holding is not a financial risk criterion.

\begin{figure}[tb]
\centering
\definecolor{hedge0}{HTML}{0072B2}
\definecolor{hedge1}{HTML}{D55E00}
\definecolor{hedge2}{HTML}{009E73}
\begin{tikzpicture}
\begin{groupplot}[group style={group size=2 by 1,horizontal sep=13mm},
width=0.45\textwidth,height=5.6cm,xmin=0,xmax=4,grid=major,
grid style={gray!25},tick label style={font=\scriptsize},
label style={font=\small},title style={font=\small},
xlabel={Budget/premium $\kappa$}]
\nextgroupplot[title={(a) Risk under funding limits},ylabel={$v(\kappa V_0)/L(0)$},
ymin=0,ymax=1.05,legend style={font=\scriptsize,draw=none,fill=none,at={(.98,.98)},anchor=north east}]
\addplot[color=hedge0,thick] coordinates {(0,1) (0.033333333,0.9258812958) (0.066666667,0.8555567) (0.1,0.7890262124) (0.13333333,0.726289833) (0.16666667,0.667347562) (0.2,0.6121993992) (0.23333333,0.5608453447) (0.26666667,0.5132853984) (0.3,0.4695195605) (0.33333333,0.4295478308) (0.36666667,0.3933702094) (0.4,0.3609866963) (0.43333333,0.3323972914) (0.46666667,0.3076019948) (0.5,0.2866008065) (0.53333333,0.2693937265) (0.56666667,0.2559807548) (0.6,0.2463618913) (0.63333333,0.2405371361) (0.66666667,0.2384263004) (0.7,0.2375189206) (0.73333333,0.2366219423) (0.76666667,0.2357353654) (0.8,0.2348591899) (0.83333333,0.2339934158) (0.86666667,0.2331380432) (0.9,0.2322930719) (0.93333333,0.2314585021) (0.96666667,0.2306343337) (1,0.2298205667) (1.0333333,0.2290172011) (1.0666667,0.2282242369) (1.1,0.2274416741) (1.1333333,0.2266695128) (1.1666667,0.2259077529) (1.2,0.2251563943) (1.2333333,0.2244154372) (1.2666667,0.2236848815) (1.3,0.2229647273) (1.3333333,0.2222549744) (1.3666667,0.2215556229) (1.4,0.2208666729) (1.4333333,0.2201881243) (1.4666667,0.2195199771) (1.5,0.2188622313) (1.5333333,0.2182148869) (1.5666667,0.217577944) (1.6,0.2169514024) (1.6333333,0.2163352623) (1.6666667,0.2157295235) (1.7,0.2151341862) (1.7333333,0.2145492503) (1.7666667,0.2139747159) (1.8,0.2134105828) (1.8333333,0.2128568512) (1.8666667,0.2123135209) (1.9,0.2117805921) (1.9333333,0.2112580647) (1.9666667,0.2107459387) (2,0.2102442141) (2.0333333,0.209752891) (2.0666667,0.2092719692) (2.1,0.2088014489) (2.1333333,0.2083413299) (2.1666667,0.2078916124) (2.2,0.2074522963) (2.2333333,0.2070233817) (2.2666667,0.2066048684) (2.3,0.2061967565) (2.3333333,0.2057990461) (2.3666667,0.2054117371) (2.4,0.2050348295) (2.4333333,0.2046683233) (2.4666667,0.2043122185) (2.5,0.2039665151) (2.5333333,0.2036312132) (2.5666667,0.2033063126) (2.6,0.2029918135) (2.6333333,0.2026877158) (2.6666667,0.2023940195) (2.7,0.2021107246) (2.7333333,0.2018378312) (2.7666667,0.2015753391) (2.8,0.2013232485) (2.8333333,0.2010815592) (2.8666667,0.2008502714) (2.9,0.200629385) (2.9333333,0.2004189001) (2.9666667,0.2002188165) (3,0.2000291343) (3.0333333,0.1998498536) (3.0666667,0.1996809743) (3.1,0.1995224964) (3.1333333,0.1993744199) (3.1666667,0.1992367448) (3.2,0.1991094711) (3.2333333,0.1989925989) (3.2666667,0.198886128) (3.3,0.1987900586) (3.3333333,0.1987043906) (3.3666667,0.198629124) (3.4,0.1985642588) (3.4333333,0.198509795) (3.4666667,0.1984657327) (3.5,0.1984320717) (3.5333333,0.1984088122) (3.5666667,0.1983959541) (3.6,0.1983931355) (3.6333333,0.1983931355) (3.6666667,0.1983931355) (3.7,0.1983931355) (3.7333333,0.1983931355) (3.7666667,0.1983931355) (3.8,0.1983931355) (3.8333333,0.1983931355) (3.8666667,0.1983931355) (3.9,0.1983931355) (3.9333333,0.1983931355) (3.9666667,0.1983931355) (4,0.1983931355)};
\addlegendentry{Lower kink}
\addplot[color=hedge0,dashed,forget plot] coordinates {(0,0.1983931355)(4,0.1983931355)};
\addplot[color=hedge1,thick] coordinates {(0,1) (0.033333333,0.9219130792) (0.066666667,0.8484661204) (0.1,0.7796591234) (0.13333333,0.7154920884) (0.16666667,0.6559650152) (0.2,0.601077904) (0.23333333,0.5508307547) (0.26666667,0.5052235673) (0.3,0.4642563417) (0.33333333,0.4279290781) (0.36666667,0.3962417764) (0.4,0.3691944366) (0.43333333,0.3467870587) (0.46666667,0.3290196427) (0.5,0.3158921887) (0.53333333,0.3074046965) (0.56666667,0.3035571662) (0.6,0.302261248) (0.63333333,0.3009837864) (0.66666667,0.2997190126) (0.7,0.2984669266) (0.73333333,0.2972275285) (0.76666667,0.2960008182) (0.8,0.2947867957) (0.83333333,0.2935854611) (0.86666667,0.2923968143) (0.9,0.2912208554) (0.93333333,0.2900575843) (0.96666667,0.2889070011) (1,0.2877691056) (1.0333333,0.2866438981) (1.0666667,0.2855313783) (1.1,0.2844315464) (1.1333333,0.2833444023) (1.1666667,0.2822699461) (1.2,0.2812081777) (1.2333333,0.2801590972) (1.2666667,0.2791227045) (1.3,0.2780989996) (1.3333333,0.2770879826) (1.3666667,0.2760896534) (1.4,0.275104012) (1.4333333,0.2741310585) (1.4666667,0.2731707928) (1.5,0.272223215) (1.5333333,0.271288325) (1.5666667,0.2703661228) (1.6,0.2694566085) (1.6333333,0.268559782) (1.6666667,0.2676756434) (1.7,0.2668041926) (1.7333333,0.2659454296) (1.7666667,0.2650993545) (1.8,0.2642659672) (1.8333333,0.2634452677) (1.8666667,0.2626372561) (1.9,0.2618419323) (1.9333333,0.2610592964) (1.9666667,0.2602893483) (2,0.259532088) (2.0333333,0.2587875156) (2.0666667,0.258055631) (2.1,0.2573364343) (2.1333333,0.2566299254) (2.1666667,0.2559361043) (2.2,0.2552549711) (2.2333333,0.2545865257) (2.2666667,0.2539307682) (2.3,0.2532876985) (2.3333333,0.2526573166) (2.3666667,0.2520396226) (2.4,0.2514346164) (2.4333333,0.250842298) (2.4666667,0.2502626675) (2.5,0.2496957248) (2.5333333,0.24914147) (2.5666667,0.248599903) (2.6,0.2480710238) (2.6333333,0.2475548325) (2.6666667,0.247051329) (2.7,0.2465605134) (2.7333333,0.2460823856) (2.7666667,0.2456169456) (2.8,0.2451641935) (2.8333333,0.2447241292) (2.8666667,0.2442967527) (2.9,0.2438820641) (2.9333333,0.2434800634) (2.9666667,0.2430907504) (3,0.2427141253) (3.0333333,0.2423501881) (3.0666667,0.2419989387) (3.1,0.2416603771) (3.1333333,0.2413345033) (3.1666667,0.2410213174) (3.2,0.2407208194) (3.2333333,0.2404330091) (3.2666667,0.2401578868) (3.3,0.2398954522) (3.3333333,0.2396457055) (3.3666667,0.2394086466) (3.4,0.2391842756) (3.4333333,0.2389725924) (3.4666667,0.2387735971) (3.5,0.2385872896) (3.5333333,0.2384136699) (3.5666667,0.238252738) (3.6,0.238104494) (3.6333333,0.2379689379) (3.6666667,0.2378460696) (3.7,0.2377358891) (3.7333333,0.2376383964) (3.7666667,0.2375535916) (3.8,0.2374814747) (3.8333333,0.2374220455) (3.8666667,0.2373753043) (3.9,0.2373412508) (3.9333333,0.2373198852) (3.9666667,0.2373112074) (4,0.2373109928)};
\addlegendentry{Negative rate/dividend}
\addplot[color=hedge1,dashed,forget plot] coordinates {(0,0.2373109928)(4,0.2373109928)};
\addplot[color=hedge2,thick] coordinates {(0,1) (0.033333333,0.930119223) (0.066666667,0.8679453333) (0.1,0.8134783307) (0.13333333,0.7667182154) (0.16666667,0.7276649874) (0.2,0.6963186465) (0.23333333,0.6726791929) (0.26666667,0.6567466265) (0.3,0.6462828058) (0.33333333,0.6362063827) (0.36666667,0.6263216275) (0.4,0.6166285402) (0.43333333,0.6071271207) (0.46666667,0.597817369) (0.5,0.5886992852) (0.53333333,0.5797728693) (0.56666667,0.5710381212) (0.6,0.5624950409) (0.63333333,0.5541436285) (0.66666667,0.5459838839) (0.7,0.5380158072) (0.73333333,0.5302573559) (0.76666667,0.5234030545) (0.8,0.5178094172) (0.83333333,0.513476444) (0.86666667,0.5104041348) (0.9,0.5085924897) (0.93333333,0.5080390114) (0.96666667,0.5080390114) (1,0.5080390114) (1.0333333,0.5080390114) (1.0666667,0.5080390114) (1.1,0.5080390114) (1.1333333,0.5080390114) (1.1666667,0.5080390114) (1.2,0.5080390114) (1.2333333,0.5080390114) (1.2666667,0.5080390114) (1.3,0.5080390114) (1.3333333,0.5080390114) (1.3666667,0.5080390114) (1.4,0.5080390114) (1.4333333,0.5080390114) (1.4666667,0.5080390114) (1.5,0.5080390114) (1.5333333,0.5080390114) (1.5666667,0.5080390114) (1.6,0.5080390114) (1.6333333,0.5080390114) (1.6666667,0.5080390114) (1.7,0.5080390114) (1.7333333,0.5080390114) (1.7666667,0.5080390114) (1.8,0.5080390114) (1.8333333,0.5080390114) (1.8666667,0.5080390114) (1.9,0.5080390114) (1.9333333,0.5080390114) (1.9666667,0.5080390114) (2,0.5080390114) (2.0333333,0.5080390114) (2.0666667,0.5080390114) (2.1,0.5080390114) (2.1333333,0.5080390114) (2.1666667,0.5080390114) (2.2,0.5080390114) (2.2333333,0.5080390114) (2.2666667,0.5080390114) (2.3,0.5080390114) (2.3333333,0.5080390114) (2.3666667,0.5080390114) (2.4,0.5080390114) (2.4333333,0.5080390114) (2.4666667,0.5080390114) (2.5,0.5080390114) (2.5333333,0.5080390114) (2.5666667,0.5080390114) (2.6,0.5080390114) (2.6333333,0.5080390114) (2.6666667,0.5080390114) (2.7,0.5080390114) (2.7333333,0.5080390114) (2.7666667,0.5080390114) (2.8,0.5080390114) (2.8333333,0.5080390114) (2.8666667,0.5080390114) (2.9,0.5080390114) (2.9333333,0.5080390114) (2.9666667,0.5080390114) (3,0.5080390114) (3.0333333,0.5080390114) (3.0666667,0.5080390114) (3.1,0.5080390114) (3.1333333,0.5080390114) (3.1666667,0.5080390114) (3.2,0.5080390114) (3.2333333,0.5080390114) (3.2666667,0.5080390114) (3.3,0.5080390114) (3.3333333,0.5080390114) (3.3666667,0.5080390114) (3.4,0.5080390114) (3.4333333,0.5080390114) (3.4666667,0.5080390114) (3.5,0.5080390114) (3.5333333,0.5080390114) (3.5666667,0.5080390114) (3.6,0.5080390114) (3.6333333,0.5080390114) (3.6666667,0.5080390114) (3.7,0.5080390114) (3.7333333,0.5080390114) (3.7666667,0.5080390114) (3.8,0.5080390114) (3.8333333,0.5080390114) (3.8666667,0.5080390114) (3.9,0.5080390114) (3.9333333,0.5080390114) (3.9666667,0.5080390114) (4,0.5080390114)};
\addlegendentry{High volatility}
\addplot[color=hedge2,dashed,forget plot] coordinates {(0,0.3763501797)(4,0.3763501797)};
\addplot[gray,dotted,forget plot] coordinates {(0,1)(4,1)};
\nextgroupplot[title={(b) Value of additional budget},ylabel={$\lambda_B V_0/L(0)$},ymin=0]
\addplot[color=hedge0,thick] coordinates {(0,2.280472749) (0.033333333,2.166649501) (0.066666667,2.052826252) (0.1,1.939003004) (0.13333333,1.825179756) (0.16666667,1.711356508) (0.2,1.597533259) (0.23333333,1.483710011) (0.26666667,1.369886763) (0.3,1.256063515) (0.33333333,1.142240266) (0.36666667,1.028417018) (0.4,0.9145937698) (0.43333333,0.8007705215) (0.46666667,0.6869472733) (0.5,0.573124025) (0.53333333,0.4593007767) (0.56666667,0.3454775285) (0.6,0.2316542802) (0.63333333,0.117831032) (0.66666667,0.02737741319) (0.7,0.02706537076) (0.73333333,0.02675332834) (0.76666667,0.02644128591) (0.8,0.02612924349) (0.83333333,0.02581720106) (0.86666667,0.02550515864) (0.9,0.02519311621) (0.93333333,0.02488107378) (0.96666667,0.02456903136) (1,0.02425698893) (1.0333333,0.02394494651) (1.0666667,0.02363290408) (1.1,0.02332086166) (1.1333333,0.02300881923) (1.1666667,0.02269677681) (1.2,0.02238473438) (1.2333333,0.02207269195) (1.2666667,0.02176064953) (1.3,0.0214486071) (1.3333333,0.02113656468) (1.3666667,0.02082452225) (1.4,0.02051247983) (1.4333333,0.0202004374) (1.4666667,0.01988839497) (1.5,0.01957635255) (1.5333333,0.01926431012) (1.5666667,0.0189522677) (1.6,0.01864022527) (1.6333333,0.01832818285) (1.6666667,0.01801614042) (1.7,0.01770409799) (1.7333333,0.01739205557) (1.7666667,0.01708001314) (1.8,0.01676797072) (1.8333333,0.01645592829) (1.8666667,0.01614388587) (1.9,0.01583184344) (1.9333333,0.01551980101) (1.9666667,0.01520775859) (2,0.01489571616) (2.0333333,0.01458367374) (2.0666667,0.01427163131) (2.1,0.01395958889) (2.1333333,0.01364754646) (2.1666667,0.01333550403) (2.2,0.01302346161) (2.2333333,0.01271141918) (2.2666667,0.01239937676) (2.3,0.01208733433) (2.3333333,0.01177529191) (2.3666667,0.01146324948) (2.4,0.01115120706) (2.4333333,0.01083916463) (2.4666667,0.0105271222) (2.5,0.01021507978) (2.5333333,0.009903037352) (2.5666667,0.009590994927) (2.6,0.009278952501) (2.6333333,0.008966910075) (2.6666667,0.00865486765) (2.7,0.008342825224) (2.7333333,0.008030782798) (2.7666667,0.007718740373) (2.8,0.007406697947) (2.8333333,0.007094655521) (2.8666667,0.006782613096) (2.9,0.00647057067) (2.9333333,0.006158528244) (2.9666667,0.005846485819) (3,0.005534443393) (3.0333333,0.005222400967) (3.0666667,0.004910358542) (3.1,0.004598316116) (3.1333333,0.00428627369) (3.1666667,0.003974231264) (3.2,0.003662188839) (3.2333333,0.003350146413) (3.2666667,0.003038103987) (3.3,0.002726061562) (3.3333333,0.002414019136) (3.3666667,0.00210197671) (3.4,0.001789934285) (3.4333333,0.001477891859) (3.4666667,0.001165849433) (3.5,0.0008538070077) (3.5333333,0.000541764582) (3.5666667,0.0002297221564) (3.6,0) (3.6333333,0) (3.6666667,0) (3.7,0) (3.7333333,0) (3.7666667,0) (3.8,0) (3.8333333,0) (3.8666667,0) (3.9,0) (3.9333333,0) (3.9666667,0) (4,0)};
\addplot[color=hedge1,thick] coordinates {(0,2.412207052) (0.033333333,2.273008194) (0.066666667,2.133809337) (0.1,1.99461048) (0.13333333,1.855411623) (0.16666667,1.716212766) (0.2,1.577013908) (0.23333333,1.437815051) (0.26666667,1.298616194) (0.3,1.159417337) (0.33333333,1.02021848) (0.36666667,0.8810196225) (0.4,0.7418207653) (0.43333333,0.6026219082) (0.46666667,0.463423051) (0.5,0.3242241938) (0.53333333,0.1850253366) (0.56666667,0.04582647942) (0.6,0.03851416697) (0.63333333,0.0381335318) (0.66666667,0.03775289664) (0.7,0.03737226148) (0.73333333,0.03699162632) (0.76666667,0.03661099115) (0.8,0.03623035599) (0.83333333,0.03584972083) (0.86666667,0.03546908567) (0.9,0.03508845051) (0.93333333,0.03470781534) (0.96666667,0.03432718018) (1,0.03394654502) (1.0333333,0.03356590986) (1.0666667,0.03318527469) (1.1,0.03280463953) (1.1333333,0.03242400437) (1.1666667,0.03204336921) (1.2,0.03166273404) (1.2333333,0.03128209888) (1.2666667,0.03090146372) (1.3,0.03052082856) (1.3333333,0.0301401934) (1.3666667,0.02975955823) (1.4,0.02937892307) (1.4333333,0.02899828791) (1.4666667,0.02861765275) (1.5,0.02823701758) (1.5333333,0.02785638242) (1.5666667,0.02747574726) (1.6,0.0270951121) (1.6333333,0.02671447694) (1.6666667,0.02633384177) (1.7,0.02595320661) (1.7333333,0.02557257145) (1.7666667,0.02519193629) (1.8,0.02481130112) (1.8333333,0.02443066596) (1.8666667,0.0240500308) (1.9,0.02366939564) (1.9333333,0.02328876047) (1.9666667,0.02290812531) (2,0.02252749015) (2.0333333,0.02214685499) (2.0666667,0.02176621983) (2.1,0.02138558466) (2.1333333,0.0210049495) (2.1666667,0.02062431434) (2.2,0.02024367918) (2.2333333,0.01986304401) (2.2666667,0.01948240885) (2.3,0.01910177369) (2.3333333,0.01872113853) (2.3666667,0.01834050336) (2.4,0.0179598682) (2.4333333,0.01757923304) (2.4666667,0.01719859788) (2.5,0.01681796272) (2.5333333,0.01643732755) (2.5666667,0.01605669239) (2.6,0.01567605723) (2.6333333,0.01529542207) (2.6666667,0.0149147869) (2.7,0.01453415174) (2.7333333,0.01415351658) (2.7666667,0.01377288142) (2.8,0.01339224626) (2.8333333,0.01301161109) (2.8666667,0.01263097593) (2.9,0.01225034077) (2.9333333,0.01186970561) (2.9666667,0.01148907044) (3,0.01110843528) (3.0333333,0.01072780012) (3.0666667,0.01034716496) (3.1,0.009966529794) (3.1333333,0.009585894632) (3.1666667,0.00920525947) (3.2,0.008824624308) (3.2333333,0.008443989145) (3.2666667,0.008063353983) (3.3,0.007682718821) (3.3333333,0.007302083658) (3.3666667,0.006921448496) (3.4,0.006540813334) (3.4333333,0.006160178172) (3.4666667,0.005779543009) (3.5,0.005398907847) (3.5333333,0.005018272685) (3.5666667,0.004637637522) (3.6,0.00425700236) (3.6333333,0.003876367198) (3.6666667,0.003495732036) (3.7,0.003115096873) (3.7333333,0.002734461711) (3.7666667,0.002353826549) (3.8,0.001973191386) (3.8333333,0.001592556224) (3.8666667,0.001211921062) (3.9,0.0008312858995) (3.9333333,0.0004506507373) (3.9666667,7.001557496e-05) (4,0)};
\addplot[color=hedge2,thick] coordinates {(0,2.212026618) (0.033333333,1.980820001) (0.066666667,1.749613384) (0.1,1.518406768) (0.13333333,1.287200151) (0.16666667,1.055993534) (0.2,0.8247869173) (0.23333333,0.5935803005) (0.26666667,0.3623736838) (0.3,0.3051677094) (0.33333333,0.2994176739) (0.36666667,0.2936676384) (0.4,0.2879176029) (0.43333333,0.2821675674) (0.46666667,0.276417532) (0.5,0.2706674965) (0.53333333,0.264917461) (0.56666667,0.2591674255) (0.6,0.25341739) (0.63333333,0.2476673545) (0.66666667,0.241917319) (0.7,0.2361672835) (0.73333333,0.2245390032) (0.76666667,0.1867190807) (0.8,0.1488991581) (0.83333333,0.1110792356) (0.86666667,0.07325931308) (0.9,0.03543939055) (0.93333333,0) (0.96666667,0) (1,0) (1.0333333,0) (1.0666667,0) (1.1,0) (1.1333333,0) (1.1666667,0) (1.2,0) (1.2333333,0) (1.2666667,0) (1.3,0) (1.3333333,0) (1.3666667,0) (1.4,0) (1.4333333,0) (1.4666667,0) (1.5,0) (1.5333333,0) (1.5666667,0) (1.6,0) (1.6333333,0) (1.6666667,0) (1.7,0) (1.7333333,0) (1.7666667,0) (1.8,0) (1.8333333,0) (1.8666667,0) (1.9,0) (1.9333333,0) (1.9666667,0) (2,0) (2.0333333,0) (2.0666667,0) (2.1,0) (2.1333333,0) (2.1666667,0) (2.2,0) (2.2333333,0) (2.2666667,0) (2.3,0) (2.3333333,0) (2.3666667,0) (2.4,0) (2.4333333,0) (2.4666667,0) (2.5,0) (2.5333333,0) (2.5666667,0) (2.6,0) (2.6333333,0) (2.6666667,0) (2.7,0) (2.7333333,0) (2.7666667,0) (2.8,0) (2.8333333,0) (2.8666667,0) (2.9,0) (2.9333333,0) (2.9666667,0) (3,0) (3.0333333,0) (3.0666667,0) (3.1,0) (3.1333333,0) (3.1666667,0) (3.2,0) (3.2333333,0) (3.2666667,0) (3.3,0) (3.3333333,0) (3.3666667,0) (3.4,0) (3.4333333,0) (3.4666667,0) (3.5,0) (3.5333333,0) (3.5666667,0) (3.6,0) (3.6333333,0) (3.6666667,0) (3.7,0) (3.7333333,0) (3.7666667,0) (3.8,0) (3.8333333,0) (3.8666667,0) (3.9,0) (3.9333333,0) (3.9666667,0) (4,0)};
\end{groupplot}
\end{tikzpicture}
\caption{Funding and residual risk for three registered contractual scenarios at $\mu=0.06$. Solid curves use the constrained mean-square hedge; dashed horizontal levels give each two-instrument unrestricted relaxation. The dotted level is the cash-only baseline. The right panel shows the normalised budget multiplier, equal to the marginal reduction in normalised risk on regular active-set intervals. At zero I use the proved right derivative. A supplemental 121-point grid illustrates curvature and saturation; all seven contracts remain in the primary comparison. These are synthetic floating-point diagnostics.}
\label{fig:h-budget-frontier}
\end{figure}
Figure~\ref{fig:h-budget-frontier} links the funding constraint to
both residual risk and its local marginal value. I use two-dimensional
curves because each panel varies one budget parameter. The supplemental
grid adds resolution to the function, not further selected hedge trials.
The primary tables retain every registered contract and physical law.

I also evaluate all holdings on 65,536 independent normal paths per
contract/physical-law regime, using seed 674193 plus
$1009(3i+j)$ for zero-based contract and drift indices.
The same paths compare the methods within a regime, yielding 420
Monte Carlo rows. No path is used to fit a hedge and no raw paths
are retained. The reported standard errors concern sample means of
squared residuals and are asymptotic diagnostics. All residual fourth
moments are finite under the stipulated lognormal law, supporting
the usual independent-sample convergence assumptions; these standard
errors are not finite-sample coverage certificates.
The maximum absolute standardised discrepancy is $2.21075$.
Two rows exceed $1.96$, but they use identical holdings and the same
paths, so they are not independent observations or simultaneous
coverage evidence.

Three negative expanded-quadratic risks and nine heuristic-floor
constrained rows remain in the archive. The floor is
$128\epsilon_{\mathrm{mach}}$ times the sum of absolute quadratic
terms; it is a diagnostic choice, not a proved roundoff bound.
For the short-dated case the residual-square calculation gives values
around $10^{-30}$--$10^{-28}$ for near-unit-call holdings, while
moment subtraction can report zero or a negative number.
The true two-instrument optimum is strictly positive by the proof.
Numerical positivity alone does not resolve its magnitude: tiny
coefficient/funding errors and rare unsampled tails can dominate
these computed risks.

\section{Interpretation and reproducibility}

The financial question is how much discounted terminal error remains
when a fixed fair premium must fund a restricted static hedge.
The exact argument separates initial valuation under $Q$, physical
loss under $P$, funding feasibility and numerical error. Its useful
project-specific conclusions are the loss distinction, the exact
constraint penalty and the interpretable budget frontier. These are
checked specialisations of established hedging and convex optimisation,
with reproducible adverse comparisons.

The engine is \path{gbm_hedging.py}. Its runner,
\path{run_hedging_study.py}, reads
\path{hedging_protocol.json} and writes a fresh result
directory. Python 3.13, NumPy 2.3.5 and SciPy 1.16.3 reproduce the
registered computations. Ten tests compare independent moment
representations, active-set and general-purpose solutions, budget
subgradients, dividend funding and cancellation cases.
The source ledger records publication status, actual consulted versions
and reading scope. Full derivations follow in the appendix.
The existing pricing/calibration companion and main quantum paper
retain their previous evidence.

\appendix
\section{Moment, optimality and sensitivity proofs}
\label{app:h-proofs}

\subsection{Discounted financial accounting and exact moments}

A fixed unit of the total-return stock has terminal value $e^{qT}X$;
continuous dividends are reinvested, making the share count evolve
mechanically as $e^{qt}$. A call pays $(X-K)^+$ and the
bank balance accumulates at $r$. Thus discounted portfolio wealth is
$a\theta_sX+D\theta_c(X-K)^++V_0-p^\top\theta
=V_0+\theta^\top g$. For the constrained long-only positions,
the available securities have non-negative prices. Discounted wealth
is bounded below by the fixed initial bank balance, itself at least
$V_0-B$. This proves admissibility on the finite horizon even when
borrowing occurs. The unrestricted $L^2$ relaxation need not satisfy
that condition. Expanding the squared residual gives \eqref{eq:h-risk}.
All moments are finite because lognormal variables have every positive
integer moment.

Write $b=\log S+(\mu-q-\sigma^2/2)T$, $v=\sigma\sqrt T$ and
$t_j=\mathbb E_PX^j=\exp(jb+j^2v^2/2)$. For $j=0,1,2$, define
\begin{align}
 T_j(l,u)&=\mathbb E_P[X^j{\bf1}_{\{l<X\leq u\}}]\notag\\
 &=t_j\left[
 \Phi\!\left(\frac{\log u-b-jv^2}{v}\right)-
 \Phi\!\left(\frac{\log l-b-jv^2}{v}\right)\right].
 \label{eq:h-truncated-moment}
\end{align}
The conventions $\log0=-\infty$ and $\log\infty=\infty$ cover tails.
To derive this formula, integrate $e^{j(b+vz)}\varphi(z)$ over the
corresponding normal interval. Completing the square,
$jvz-z^2/2=-(z-jv)^2/2+j^2v^2/2$, gives the displayed Gaussian
probability. Interval endpoints have zero probability.

Let $F=\min\{(X-K)^+,H\}$ and $A=(X-K)^+$. Put
$f_0=\mathbb EF$, $f_1=\mathbb EXF$, $f_2=\mathbb EF^2$ and
$a_0=\mathbb EA$, $a_1=\mathbb EXA$, $a_2=\mathbb EA^2$.
Splitting at $K,J$ gives
\begin{align}
 f_0&=T_1(K,J)-KT_0(K,J)+HT_0(J,\infty),\notag\\
 f_1&=T_2(K,J)-KT_1(K,J)+HT_1(J,\infty),\notag\\
 f_2&=T_2(K,J)-2KT_1(K,J)+K^2T_0(K,J)+H^2T_0(J,\infty),
 \label{eq:h-capped-moments}\\
 a_0&=T_1(K,\infty)-KT_0(K,\infty),\notag\\
 a_1&=T_2(K,\infty)-KT_1(K,\infty),\notag\\
 a_2&=T_2(K,\infty)-2KT_1(K,\infty)+K^2T_0(K,\infty),
 \label{eq:h-call-moments}\\
 z&=\mathbb EFA=T_2(K,J)-2KT_1(K,J)+K^2T_0(K,J)\notag\\
 &\hspace{14mm}+H\{T_1(J,\infty)-KT_0(J,\infty)\}.
 \label{eq:h-payoff-cross-moment}
\end{align}
Expansion of the instrument gains now supplies
\begin{align}
 m&=(at_1-S,\ Da_0-C_K)^\top,\notag\\
 M_{ss}&=a^2t_2-2aSt_1+S^2,\notag\\
 M_{cc}&=D^2a_2-2DC_Ka_0+C_K^2,\notag\\
 M_{sc}&=aDa_1-aC_Kt_1-SDa_0+SC_K,\label{eq:h-gain-moments}\\
 h_s&=aDf_1-SDf_0-V_0(at_1-S),\notag\\
 h_c&=D^2z-DC_Kf_0-V_0(Da_0-C_K),\notag\\
 c&=D^2f_2-2DV_0f_0+V_0^2,\qquad e=Df_0-V_0.
 \label{eq:h-loss-moments}
\end{align}

There is a separate positive-integral representation. On the successive
regions $(0,K)$, $(K,J)$ and $(J,\infty)$, write $R=A_iX+B_i$.
The coefficients are
\begin{align}
 (A_1,A_2,A_3)&=(-a\theta_s,\ D(1-\theta_c)-a\theta_s,\
 -D\theta_c-a\theta_s),\notag\\
 B_1&=-V_0+S\theta_s+C_K\theta_c,\notag\\
 B_2&=-DK(1-\theta_c)-V_0+S\theta_s+C_K\theta_c,\notag\\
 B_3&=DH+DK\theta_c-V_0+S\theta_s+C_K\theta_c.
 \label{eq:h-affine-residual}
\end{align}
Integrating $\sum_i{\bf1}_i(A_iX+B_i)^2$ against the positive
lognormal density gives the same risk. This mathematical identity
does not guarantee that either expanded floating-point calculation
is free from cancellation.

\subsection{Positive definiteness, existence and the risk penalty}

The lognormal density is strictly positive on $(0,\infty)$.
If $u^\top Mu=\mathbb E(u^\top g)^2=0$, continuity makes
$u^\top g(x)=0$ on that whole interval. Below $K$, its affine slope
is $au_s$, so $u_s=0$, and its constant is $-u_cC_K$, so
$u_c=0$ because $C_K>0$. Thus $M\succ0$.
For $G=\operatorname{Cov}(g)$, zero quadratic variance would instead
make $u^\top g(x)$ constant. The below-$K$ slope again gives
$u_s=0$, while the above-$K$ slope $Du_c$ gives $u_c=0$.
Hence $G\succ0$ as well.

The feasible set is closed, bounded and contains zero. Continuity
therefore gives existence, and the positive-definite Hessian $2M$
gives strict convexity and uniqueness. Completing the quadratic yields
\begin{equation}
 L(\theta)=c-h^\top M^{-1}h+
 (\theta-\theta_M)^\top M(\theta-\theta_M),\qquad
 \theta_M=M^{-1}h.
 \label{eq:h-square-completion}
\end{equation}
This proves the constraint penalty stated in the main text.
Nevertheless the unrestricted minimum is strictly positive. If any
finite $\theta$ made $L=0$, continuity and positive density would
make all three affine residuals identically zero. Their slopes
in \eqref{eq:h-affine-residual} force, successively, $\theta_s=0$,
$\theta_c=1$ and $\theta_c=0$, a contradiction. The unrestricted
minimum is attained at $\theta_M$, so this is a strictly positive
minimum, not merely positivity at each chosen hedge.

\subsection{KKT necessity without a constraint qualification}

Write the five inequalities as $A\theta\leq b(B)$, with rows
$-e_s^\top,-e_c^\top,e_s^\top,e_c^\top,p^\top$ and right sides
$0,0,1,1,B$. At a minimiser $\theta_*$, any direction $d$ with
$A_Id\leq0$ for the active rows $I$ remains feasible for all sufficiently
small positive steps: each inactive inequality has a positive slack.
Differentiability gives $\nabla L(\theta_*)^\top d\geq0$.
For completeness, let $K_I$ be the cone of the active row normals.
This cone is closed in $\mathbb R^2$. Indeed any conic combination
reduces to at most two independent normals: for dependent vectors,
subtract a scalar multiple of a nonzero linear dependence from the
non-negative coefficients until one coefficient vanishes, preserving
all remaining non-negativity. Iteration terminates with independence.
For a convergent sequence of conic combinations, choose a subsequence
using the same independent normals, since there are finitely many
choices. Their coefficients are continuous linear functions of the
combination and hence have non-negative limits. The limit lies in
$K_I$, proving closedness.

If $w=-\nabla L(\theta_*)$ were outside $K_I$, let $k_0$ be its
nearest point in that closed convex cone and set $d=w-k_0\ne0$.
Existence follows by minimising distance on a sufficiently large compact ball.
Projection optimality gives $d^\top(k-k_0)\leq0$ for every $k\in K_I$.
Taking $k=0,2k_0$ gives $d^\top k_0=0$, hence $d^\top k\leq0$
for every $k\in K_I$. Thus $A_Id\leq0$, but
$\nabla L(\theta_*)^\top d=-w^\top d=-\|d\|^2<0$,
contradicting the directional inequality. Therefore $w\in K_I$
and has non-negative conic coefficients. These are the active
multipliers satisfying $\nabla L+A_I^\top\lambda_I=0$; inactive
multipliers are zero. This proves the polyhedral Farkas argument here
and yields \eqref{eq:h-kkt-stationarity}--\eqref{eq:h-kkt-complementarity}.
Conversely, for any feasible $\theta$, convexity gives
$L(\theta)\geq L(\theta_*)+\nabla L(\theta_*)^\top(\theta-\theta_*)$.
Stationarity and complementary slackness make the latter inner product
non-negative, proving sufficiency. No strict feasible point is used.

At $B=0$, positivity of $p$ makes zero the only feasible point.
Set upper multipliers to zero and choose
$\lambda\geq\max\{0,2h_s/S,2h_c/C_K\}$; then lower multipliers
$\ell=\lambda p-2h$ are non-negative and all KKT equations hold.
For $B>0$, sufficiently small positive positions also demonstrate
strict feasibility, but necessity does not depend on that observation.

For the exact candidate method, an interior optimum has
$\theta=M^{-1}h$. A relative interior point of any nondegenerate face
lies on one of the five supporting lines. Writing that line as
$x_0+tv$, $v\ne0$, its restricted objective has positive second
derivative $2v^\top Mv$, and its unique stationary parameter is
\begin{equation}
 t_*=\frac{v^\top(h-Mx_0)}{v^\top Mv}.
 \label{eq:h-face-stationary}
\end{equation}
Keep this candidate only if feasible. Every vertex is an intersection
of two independent active line normals. There are at most ten line
pairs; parallel pairs have no unique intersection. These alternatives
cover an optimum on every face, including the singleton at $B=0$.
Testing at most $1+5+10=16$ candidates therefore finds the global
minimum; duplicate candidates and empty faces can be discarded.

\subsection{Why a variance control is a different financial hedge}

Expanding means and covariances gives $M=G+mm^\top$ and
$h=k+me$, hence
\begin{equation}
 L(\theta)=\operatorname{Var}(Y-\theta^\top g)
 +(e-m^\top\theta)^2.
 \label{eq:h-mean-variance-decomposition}
\end{equation}
The variance minimiser is $\theta_v=G^{-1}k$. The identity
$(G+mm^\top)G^{-1}m=(1+\zeta)m$ shows directly that
$\theta_v+G^{-1}m\eta/(1+\zeta)$ solves $M\theta=h$.
Uniqueness proves \eqref{eq:h-centred-risk-gap}; substituting the
difference into \eqref{eq:h-square-completion} gives
$\eta^2(\zeta+\zeta^2)/(1+\zeta)^2=\eta^2\zeta/(1+\zeta)$.
Under $P=Q$, fair discounted instrument gains and claim premia have
zero means, so $m=e=0$. If the measures differ, equality can still
hold when $\eta=0$; a strict discrepancy is not automatic.

\subsection{Budget value, shadow prices and local sensitivity}

Feasible-set inclusion makes $v$ non-increasing. For $B_1,B_2\geq0$,
a convex combination of their optimisers is feasible at the same
convex combination of budgets. Convexity of $L$ then gives
$v(tB_1+(1-t)B_2)\leq tv(B_1)+(1-t)v(B_2)$ for $0\leq t\leq1$.
This proves convexity of the value function.

For any optimiser at $B$ and its KKT multipliers, the box multipliers
give non-negative terms against any other box-feasible point.
The budget term gives
$L(\theta)\geq v(B)-\lambda_B(p^\top\theta-p^\top\theta_B)$.
If $\lambda_B>0$, complementary slackness gives $p^\top\theta_B=B$;
if it is zero, the term vanishes. For $\theta\in\mathcal F_{B'}$,
the displayed bound is at least $v(B)-\lambda_B(B'-B)$.
Minimising proves \eqref{eq:h-budget-shadow}. Thus $-\lambda_B$
is a subgradient on the budget domain, even at a degenerate budget.

For local differentiability, fix an active set $I$ with independent
rows, strictly positive active multipliers and strictly inactive
remaining constraints. Put $H_0=2M$. The matrix
\begin{equation}
 \begin{pmatrix}H_0&A_I^\top\\A_I&0\end{pmatrix}
 \label{eq:h-kkt-system}
\end{equation}
is invertible: a null vector satisfies $H_0x+A_I^\top y=0$,
$A_Ix=0$, so $x^\top H_0x=0$, hence $x=0$ and then $y=0$.
The KKT solution is consequently affine in $B$. Strict inequalities
persist on an open neighbourhood, making that affine solution the
unique optimum there. If $I$ is nonempty, let $\dot b_I$ be the vector
with one at the budget row when present and zero elsewhere. Then
\begin{align}
 \dot\theta&=H_0^{-1}A_I^\top
 (A_IH_0^{-1}A_I^\top)^{-1}\dot b_I,\notag\\
 v'&=-\lambda_B,\qquad v''=\dot\theta^\top H_0\dot\theta\geq0.
 \label{eq:h-local-budget-sensitivity}
\end{align}
The first identity follows by differentiating the KKT system; the
second uses $\nabla L^\top\dot\theta=-\lambda_I^\top\dot b_I$.
For an empty active set, $\theta=\theta_M$ is locally constant.
At active-set changes or dependent rows, this differentiability
argument is not asserted; the global subgradient bound still holds.

In the budget-only active regime, put $z_p=p^\top M^{-1}p>0$.
Solving stationarity and $p^\top\theta=B$ explicitly gives
\begin{align}
 \theta&=\theta_M-\frac{p^\top\theta_M-B}{z_p}M^{-1}p,\notag\\
 \lambda_B&=\frac{2(p^\top\theta_M-B)}{z_p},\notag\\
 v(B)&=L(\theta_M)+\frac{(p^\top\theta_M-B)^2}{z_p}.
 \label{eq:h-budget-only-regime}
\end{align}
These formulas apply only while the box is strictly inactive and
$\lambda_B>0$.

Finally, for sufficiently small $B>0$, write $\theta=Bz$.
The upper bounds are then inactive on
$z\geq0,\ p^\top z\leq1$, since $z_i\leq1/p_i$.
On this compact set,
$(L(Bz)-c)/B=-2h^\top z+Bz^\top Mz$ converges uniformly to
$-2h^\top z$. Its minimum therefore tends to
$-2\max\{0,h_s/S,h_c/C_K\}$, proving the right derivative at zero.
No data-dependent active-set heuristic or estimated multiplier is
needed for this exact financial sensitivity.

\subsection{A bound from feasible KKT diagnostics}

For any exactly feasible $\theta$ and non-negative multipliers
$\lambda$, put $r_*=\nabla L(\theta)+A^\top\lambda$,
$\varepsilon_c=\lambda^\top(b-A\theta)\geq0$ and
$m_*=\lambda_{\min}(M)>0$. With $\delta=\theta_*-\theta$,
quadratic expansion and feasibility give
$L(\theta_*)\geq L(\theta)+r_*^\top\delta+
m_*\|\delta\|^2-\varepsilon_c$. Completing the square and using
the optimum's variational inequality yield
\begin{equation}
 m_*\|\theta-\theta_*\|^2
 \leq L(\theta)-L(\theta_*)
 \leq \varepsilon_c+\frac{\|r_*\|^2}{4m_*}.
 \label{eq:h-kkt-risk-bound}
\end{equation}
The left inequality follows by expanding at $\theta_*$ and
using $\nabla L(\theta_*)^\top(\theta-\theta_*)\geq0$.
Floating-point feasibility, residuals and eigenvalues do not themselves
prove the exact premises or certify this bound.

\begin{thebibliography}{99}
\bibitem[Black and Scholes(1973)]{blackscholes1973}
Black, F. and Scholes, M. (1973). The pricing of options and corporate liabilities. \emph{Journal of Political Economy}, 81(3), 637--654. \url{https://doi.org/10.1086/260062}. Final journal scan consulted, printed pages 637--644.

\bibitem[Boyd and Vandenberghe(2004)]{boydvandenberghe2004}
Boyd, S. and Vandenberghe, L. (2004). \emph{Convex Optimization}. Cambridge University Press. \url{https://doi.org/10.1017/CBO9780511804441}. Author-hosted published-book PDF consulted, Sections 4.2.3, 5.2.3, 5.5.2--5.5.3 and the Farkas lemma in Section 5.8, with the normal-cone definition in Exercise 2.38(c). This source is a textbook.

\bibitem[Czichowsky and Schweizer(2012)]{czichowskyschweizer2012}
Czichowsky, C. and Schweizer, M. (2012). Convex duality in mean-variance hedging under convex trading constraints. \emph{Advances in Applied Probability}, 44(4), 1084--1112. \url{https://doi.org/10.1239/aap/1354716590}. Final publisher PDF consulted, Introduction and Section 2, printed pages 1084--1088, and Theorem 4.1/proof, page 1095.

\bibitem[Leung and Lorig(2016)]{leunglorig2016}
Leung, T. and Lorig, M. (2016). Optimal static quadratic hedging. \emph{Quantitative Finance}, 16(9), 1341--1355. \url{https://doi.org/10.1080/14697688.2016.1161229}. Published metadata/abstract verified; the consulted body is arXiv:1506.02074v2, posted 19 November 2015, Sections 2 and 3.2 and Appendix A, \url{https://arxiv.org/pdf/1506.02074v2}. The final journal body was not consulted.

\bibitem[Merton(1973)]{merton1973}
Merton, R. C. (1973). Theory of rational option pricing. \emph{The Bell Journal of Economics and Management Science}, 4(1), 141--183. \url{https://doi.org/10.2307/3003143}. Final journal PDF consulted, selected material in Sections 6--7, printed pages 162--166 and 168--171, particularly equation (44), page 170, and footnote 62, page 171.
\end{thebibliography}
\end{document}
